\ifdefined\BookMaster
\def\BookEnd{}
\else
\documentclass[a4paper,12pt,fleqn,openany,twoside,fontset=windows]{ctexbook}

\usepackage[fleqn]{amsmath}
\usepackage{amssymb,amsthm,mathrsfs}
\usepackage{geometry}
\usepackage[dvipsnames]{xcolor}
\usepackage{graphicx}
\usepackage{tikz}
\usetikzlibrary{calc}
\usepackage[most]{tcolorbox}
\usepackage{fancyhdr}
\usepackage{titlesec}
\usepackage{tocloft}
\usepackage{enumitem}
\usepackage{microtype}
\usepackage{pifont}
\usepackage{hyperref}
\usepackage{romannum}
\usepackage{mathtools}

\definecolor{bookblue}{HTML}{264653}
\definecolor{bookteal}{HTML}{4F7C82}
\definecolor{booklight}{HTML}{EDF4F3}
\definecolor{bookgray}{HTML}{5E686B}

\geometry{
    inner=2.7cm,
    outer=2.5cm,
    top=2.7cm,
    bottom=2.7cm,
    headheight=18pt,
    headsep=18pt,
    footskip=25pt
}
\setlength{\mathindent}{0pt}
\setlength{\parindent}{5pt}
\setlength{\parskip}{3pt plus 1pt minus 1pt}
\linespread{1.25}
\allowdisplaybreaks[3]
\AtBeginDocument{%
    \setlength{\parindent}{5pt}
    \setlength{\abovedisplayskip}{8pt plus 2pt minus 2pt}
    \setlength{\belowdisplayskip}{8pt plus 2pt minus 2pt}
    \setlength{\abovedisplayshortskip}{5pt plus 1pt minus 1pt}
    \setlength{\belowdisplayshortskip}{5pt plus 1pt minus 1pt}
}

\hypersetup{
    unicode=true,
    colorlinks=true,
    linkcolor=bookblue,
    citecolor=bookblue,
    urlcolor=bookteal,
    bookmarksopen=true,
    pdfauthor={Chen},
    pdftitle={数学分析习题整理}
}

\renewcommand{\chaptermark}[1]{%
    \edef\@tmp{第\thechapter 章\quad #1}%
    \expandafter\markboth\expandafter{\@tmp}{}%
}
\fancypagestyle{main}{%
    \fancyhf{}
    \fancyhead[LE]{\small\color{bookgray}\nouppercase{\leftmark}}
    \fancyhead[RO]{\small\color{bookgray}数学分析习题整理}
    \fancyfoot[C]{\color{bookblue}\thepage}
    \fancyfoot[R]{\small\bfseries Chen\;\; github.com/mathlover-1/math-analysis}
    \renewcommand{\headrulewidth}{0.4pt}
    \renewcommand{\footrulewidth}{0pt}
}
\fancypagestyle{plain}{%
    \fancyhf{}
    \fancyfoot[C]{\color{bookblue}\thepage}
    \fancyfoot[R]{\small\bfseries Chen\;\; github.com/mathlover-1/math-analysis}
    \renewcommand{\headrulewidth}{0pt}
    \renewcommand{\footrulewidth}{0pt}
}
\pagestyle{main}

\titleformat{\chapter}[display]
  {\normalfont\sffamily\bfseries\color{bookblue}}
  {\raggedleft\fontsize{46}{50}\selectfont\thechapter}
  {-18pt}
  {\titlerule[1.2pt]\vspace{10pt}\filright\Huge}
\titlespacing*{\chapter}{0pt}{-15pt}{36pt}

\renewcommand{\contentsname}{目\quad 录}
\renewcommand{\cfttoctitlefont}{\hfill\sffamily\bfseries\Huge\color{bookblue}}
\renewcommand{\cftaftertoctitle}{\hfill}
\setlength{\cftbeforetoctitleskip}{5pt}
\setlength{\cftaftertoctitleskip}{28pt}
\setlength{\cftbeforechapskip}{12pt}
\renewcommand{\cftchapfont}{\sffamily\bfseries\large\color{bookblue}}
\renewcommand{\cftchappagefont}{\sffamily\bfseries\large\color{bookblue}}
\renewcommand{\cftchapleader}{\color{bookteal}\cftdotfill{2.5}}

\newtcolorbox{ti}{
    enhanced,
    breakable,
    colback=booklight!55!white,
    colframe=bookteal!35!white,
    boxrule=0.45pt,
    boxsep=0pt,
    left=10pt,
    right=10pt,
    top=9pt,
    bottom=9pt,
    before skip=14pt,
    after skip=14pt,
    arc=2pt,
    outer arc=2pt,
    before upper={\parindent=0pt}
}

\newenvironment{booknotebody}{%
    \begin{center}
    \begin{minipage}{0.84\textwidth}
    \songti\zihao{-4}\linespread{1.45}\selectfont
    \setlength{\parindent}{2\ccwd}
    \setlength{\parskip}{0.75em}
}{%
    \end{minipage}
    \end{center}
}

\newenvironment{booknotelongbody}{%
    \begingroup
    \songti\zihao{-4}\linespread{1.45}\selectfont
    \setlength{\parindent}{2\ccwd}
    \setlength{\parskip}{0.75em}
    \begin{list}{}{%
        \setlength{\leftmargin}{0.08\textwidth}
        \setlength{\rightmargin}{0.08\textwidth}
        \setlength{\topsep}{0pt}
        \setlength{\partopsep}{0pt}
        \setlength{\parsep}{0pt}
        \setlength{\itemsep}{0pt}
    }
    \item\relax
}{%
    \end{list}
    \endgroup
}

\fancypagestyle{plainnowm}{%
    \fancyhf{}
    \fancyfoot[C]{\color{bookblue}\thepage}
    \renewcommand{\headrulewidth}{0pt}
    \renewcommand{\footrulewidth}{0pt}
}

\newcommand{\booknotetitle}[2][plain]{%
    \clearpage
    \phantomsection
    \addcontentsline{toc}{chapter}{#2}
    \markboth{#2}{}
    \thispagestyle{#1}
    \vspace*{1.2cm}
    \begin{center}
        {\sffamily\heiti\bfseries\fontsize{26}{34}\selectfont\color{bookblue}#2\par}
        \vspace{0.8em}
        {\color{bookteal}\rule{4.5cm}{0.8pt}\par}
    \end{center}
    \vspace{0.6cm}
}

\renewcommand{\proofname}{证明}
\renewcommand{\qedsymbol}{\color{bookteal}\ensuremath{\blacksquare}}

\title{数学分析习题整理}
\author{Chen}
\date{}

\makeatletter
% ==================== 旧版封面/封底设计备份 ====================
% 2026-08 重设计,旧代码用 \iffalse...\fi 封存,新版见下方。
\iffalse
\renewcommand{\maketitle}{%
    \begin{titlepage}
        \thispagestyle{empty}
        \setcounter{page}{0}
        \noindent\hspace*{-\parindent}\begin{tikzpicture}[remember picture,overlay]
            \fill[bookblue] (current page.north west) rectangle ([yshift=-4.2cm]current page.north east);
            \fill[booklight] ([yshift=-4.2cm]current page.north west) rectangle (current page.south east);
            \draw[bookteal!55,line width=0.7pt]
                ([xshift=2.4cm,yshift=4.3cm]current page.south west)
                .. controls +(2.5cm,1.8cm) and +(-2.8cm,-1.3cm) ..
                ([xshift=-2.4cm,yshift=7.2cm]current page.south east);
            \draw[bookteal!35,line width=0.7pt]
                ([xshift=2.5cm,yshift=6.2cm]current page.south west)
                .. controls +(3.0cm,-1.2cm) and +(-3.0cm,1.6cm) ..
                ([xshift=-2.3cm,yshift=3.8cm]current page.south east);
        \end{tikzpicture}
        \vspace*{5.6cm}
        \noindent\color{bookblue}\rule{\textwidth}{1.4pt}\par
        \vspace{8pt}
        {\centering\sffamily\bfseries\fontsize{32}{42}\selectfont\color{bookblue}\@title\par}
        \vspace{8pt}
        \noindent\color{bookteal}\rule{\textwidth}{0.6pt}\par
        \vspace{2.6cm}
        {\centering\Large\color{bookgray}\@author\par}
        \vfill
        {\centering\color{bookteal}\large
            $\displaystyle \lim_{n\to\infty} a_n$\qquad
            $\displaystyle \int_a^b f(x)\,\mathrm{d}x$\qquad
            $\displaystyle \nabla f(x,y)$\par
        }
        \vspace*{1.5cm}
    \end{titlepage}
}

\newcommand{\makebackcover}{%
    \begingroup
    \let\cleardoublepage\clearpage
    \begin{titlepage}
        \thispagestyle{empty}
        \noindent\hspace*{-\parindent}\begin{tikzpicture}[remember picture,overlay]
            \fill[bookblue] (current page.north west) rectangle (current page.south east);
            \fill[booklight,rounded corners=7pt]
                ([xshift=2.2cm,yshift=-3.0cm]current page.north west)
                rectangle
                ([xshift=-2.2cm,yshift=3.0cm]current page.south east);
            \fill[bookteal!80]
                (current page.north west) --
                ([xshift=5.3cm]current page.north west) --
                ([xshift=2.1cm,yshift=-4.4cm]current page.north west) --
                ([yshift=-6.2cm]current page.north west) -- cycle;
            \fill[bookteal!55]
                (current page.south east) --
                ([xshift=-5.8cm]current page.south east) --
                ([xshift=-2.4cm,yshift=4.2cm]current page.south east) --
                ([yshift=6.0cm]current page.south east) -- cycle;
            \foreach \radius in {0.75,1.15,1.55,1.95} {
                \draw[bookteal!55,line width=0.55pt]
                    ([xshift=-2.9cm,yshift=-4.4cm]current page.north east)
                    circle[radius=\radius cm];
            }
            \foreach \radius in {0.65,1.05,1.45} {
                \draw[booklight!65,line width=0.55pt]
                    ([xshift=2.4cm,yshift=3.6cm]current page.south west)
                    circle[radius=\radius cm];
            }
            \fill[bookteal!16,rounded corners=3pt]
                ([xshift=3.5cm,yshift=-7.1cm]current page.north west)
                rectangle
                ([xshift=-5.0cm,yshift=7.2cm]current page.south east);
            \fill[bookteal!30,rounded corners=3pt]
                ([xshift=4.4cm,yshift=-8.0cm]current page.north west)
                rectangle
                ([xshift=-4.1cm,yshift=8.1cm]current page.south east);
            \begin{scope}[shift={([xshift=-0.75cm]current page.center)}]
                \clip[rounded corners=3pt] (-5.25,-4.75) rectangle (5.25,4.75);
                \foreach \v in {-3.6,-3.0,...,3.6} {
                    \draw[bookteal!48,line width=0.45pt,opacity=0.72]
                        plot[domain=-3.6:3.6,samples=45,smooth,variable=\u]
                        ({\u+0.55*\v},{0.28*\u-0.35*\v+0.11*(\u*\u-\v*\v)});
                }
                \foreach \u in {-3.6,-3.0,...,3.6} {
                    \draw[bookblue!42,line width=0.42pt,opacity=0.62]
                        plot[domain=-3.6:3.6,samples=45,smooth,variable=\v]
                        ({\u+0.55*\v},{0.28*\u-0.35*\v+0.11*(\u*\u-\v*\v)});
                }
                \draw[bookteal!70,line width=0.75pt,opacity=0.82]
                    plot[domain=-3.6:3.6,samples=55,smooth,variable=\u]
                    ({\u},{0.28*\u+0.11*\u*\u});
                \draw[bookblue!60,line width=0.7pt,opacity=0.75]
                    plot[domain=-3.6:3.6,samples=55,smooth,variable=\v]
                    ({0.55*\v},{-0.35*\v-0.11*\v*\v});
            \end{scope}
            \draw[bookteal!45,line width=0.65pt]
                ([xshift=3.0cm,yshift=-8.2cm]current page.north west)
                .. controls +(2.2cm,2.0cm) and +(-2.0cm,-2.1cm) ..
                ([xshift=-3.0cm,yshift=9.0cm]current page.south east);
            \draw[bookteal!25,line width=0.55pt]
                ([xshift=3.1cm,yshift=8.0cm]current page.south west)
                .. controls +(2.7cm,-1.7cm) and +(-2.6cm,1.8cm) ..
                ([xshift=-3.1cm,yshift=-8.7cm]current page.north east);
            \node[anchor=west,text=bookteal!72,scale=0.78] at
                ([xshift=2.8cm,yshift=-4.0cm]current page.north west) {%
                $\displaystyle e^{\mathrm{i}\pi}+1=0$};
            \node[anchor=east,text=bookblue!62,scale=0.67] at
                ([xshift=-2.7cm,yshift=4.0cm]current page.south east) {%
                $\displaystyle
                f(x)=\sum_{n=0}^{\infty}\frac{f^{(n)}(a)}{n!}(x-a)^n$};
        \end{tikzpicture}
        \mbox{}
    \end{titlepage}
    \endgroup
}
\fi

% ==================== 新版封面(2026-08 重设计) ====================
% 设计要点:顶部深绿纵向渐变+波浪下缘+两层半透明过渡波(解决深浅绿生硬过渡);
% 深色区螺旋线/同心圆母题;标题+英文副题+作者居中;
% 中部 Fourier 部分和逼近方波曲线;底部浅色波纹(无文字)。
\renewcommand{\maketitle}{%
    \begin{titlepage}
        \thispagestyle{empty}
        \setcounter{page}{0}
        \noindent\hspace*{-\parindent}\begin{tikzpicture}[remember picture,overlay]
            % ---------- 底色 ----------
            \fill[booklight] (current page.north west) rectangle (current page.south east);

            % ---------- 顶部深色区:纵向渐变 + 波浪下缘 ----------
            \shade[top color=bookblue, bottom color=bookteal!55!bookblue]
                (current page.north west) -- (current page.north east) --
                ([yshift=-8.6cm]current page.north east)
                .. controls ([xshift=-5.2cm,yshift=-10.6cm]current page.north east)
                         and ([xshift=6.0cm,yshift=-8.0cm]current page.north west) ..
                ([yshift=-10.0cm]current page.north west) -- cycle;

            % ---------- 过渡层 1:半透明波浪 ----------
            \fill[bookteal!60!bookblue,opacity=0.42]
                ([yshift=-8.6cm]current page.north east)
                .. controls ([xshift=-5.2cm,yshift=-10.6cm]current page.north east)
                         and ([xshift=6.0cm,yshift=-8.0cm]current page.north west) ..
                ([yshift=-10.0cm]current page.north west) --
                ([yshift=-12.6cm]current page.north west)
                .. controls ([xshift=6.5cm,yshift=-11.0cm]current page.north west)
                         and ([xshift=-5.5cm,yshift=-13.2cm]current page.north east) ..
                ([yshift=-11.4cm]current page.north east) -- cycle;

            % ---------- 过渡层 2:更浅的波浪 ----------
            \fill[bookteal!35,opacity=0.35]
                ([yshift=-10.6cm]current page.north east)
                .. controls ([xshift=-6.0cm,yshift=-12.4cm]current page.north east)
                         and ([xshift=5.0cm,yshift=-9.8cm]current page.north west) ..
                ([yshift=-11.8cm]current page.north west) --
                ([yshift=-13.8cm]current page.north west)
                .. controls ([xshift=5.5cm,yshift=-12.6cm]current page.north west)
                         and ([xshift=-6.5cm,yshift=-14.6cm]current page.north east) ..
                ([yshift=-12.8cm]current page.north east) -- cycle;

            % ---------- 深色区装饰:右上螺旋线 ----------
            \begin{scope}[shift={([xshift=-3.4cm,yshift=-4.3cm]current page.north east)}]
                \draw[booklight!45!bookteal,opacity=0.55,line width=0.6pt]
                    plot[domain=0:1440,samples=220,smooth,variable=\t]
                    ({0.0028*\t*cos(\t)},{0.0028*\t*sin(\t)});
            \end{scope}

            % ---------- 深色区装饰:左下同心圆 ----------
            \foreach \r in {0.7,1.25,1.8,2.35,2.9}{
                \draw[booklight!35!bookteal,opacity=0.5,line width=0.5pt]
                    ([xshift=0.4cm,yshift=-7.9cm]current page.north west) circle[radius=\r cm];
            }

            % ---------- 主标题 ----------
            \node[anchor=center,text=booklight] at
                ([yshift=-4.5cm]current page.north)
                {\sffamily\bfseries\fontsize{40}{46}\selectfont \@title};

            % ---------- 分隔短线 ----------
            \draw[booklight!65,line width=1.2pt]
                ([xshift=-2.6cm,yshift=-5.75cm]current page.north) --
                ([xshift=2.6cm,yshift=-5.75cm]current page.north);

            % ---------- 英文副题 ----------
            \node[anchor=center,text=booklight!70!bookteal] at
                ([yshift=-6.55cm]current page.north)
                {\itshape\large Mathematical Analysis \textperiodcentered\
                 A Collection of Exercises};

            % ---------- 作者(紧随标题下方) ----------
            \node[anchor=center,text=booklight!88] at
                ([yshift=-8.2cm]current page.north)
                {\sffamily\large \@author\quad 整理};

            % ---------- 中部:Fourier 部分和逼近方波 ----------
            \begin{scope}[shift={([xshift=-5.2cm,yshift=-18.1cm]current page.north)}]
                % 坐标轴
                \draw[bookgray!55,line width=0.5pt,->] (-0.5,0) -- (10.9,0);
                \draw[bookgray!55,line width=0.5pt,->] (0,-1.75) -- (0,1.95);
                % 目标方波(虚线)
                \draw[bookgray!75,line width=0.6pt,dashed]
                    (0,0.92) -- (5,0.92) (5,-0.92) -- (10,-0.92);
                % N=1
                \draw[bookteal!45,line width=0.7pt]
                    plot[domain=0:360,samples=90,smooth,variable=\x]
                    ({\x/36},{1.18*sin(\x)});
                % N=3
                \draw[bookteal!80,line width=0.8pt]
                    plot[domain=0:360,samples=120,smooth,variable=\x]
                    ({\x/36},{1.18*(sin(\x)+sin(3*\x)/3)});
                % N=7
                \draw[bookblue,line width=0.95pt]
                    plot[domain=0:360,samples=150,smooth,variable=\x]
                    ({\x/36},{1.18*(sin(\x)+sin(3*\x)/3+sin(5*\x)/5+sin(7*\x)/7)});
                % 刻度标注
                \node[anchor=north east,text=bookgray!80,scale=0.62] at (-0.12,0) {$0$};
                \node[anchor=north,text=bookgray!80,scale=0.62] at (5,0) {$\pi$};
                \node[anchor=north,text=bookgray!80,scale=0.62] at (10,0) {$2\pi$};
            \end{scope}
            % 图注
            \node[anchor=center,text=bookgray,scale=0.8] at
                ([yshift=-20.5cm]current page.north)
                {$S_N(x)=\displaystyle\sum_{k=1}^{N}\frac{\sin\bigl((2k-1)x\bigr)}{2k-1}$
                 \ 逐步逼近方波};

            % ---------- 底部浅色波纹(无文字,与顶部波浪呼应) ----------
            \fill[bookteal!22,opacity=0.5]
                ([yshift=2.4cm]current page.south west)
                .. controls ([xshift=6.0cm,yshift=3.9cm]current page.south west)
                         and ([xshift=-6.5cm,yshift=1.0cm]current page.south east) ..
                ([yshift=2.9cm]current page.south east) --
                (current page.south east) -- (current page.south west) -- cycle;
            \fill[bookteal!40!bookblue,opacity=0.28]
                ([yshift=1.3cm]current page.south west)
                .. controls ([xshift=5.5cm,yshift=2.7cm]current page.south west)
                         and ([xshift=-5.5cm,yshift=0.3cm]current page.south east) ..
                ([yshift=1.9cm]current page.south east) --
                (current page.south east) -- (current page.south west) -- cycle;
        \end{tikzpicture}
        \mbox{}
    \end{titlepage}
}

% ==================== 新版封底(纯背景,2026-08 重设计) ====================
% 设计要点:对角渐变底+角部半透明叠层;中下部外摆线(玫瑰形)为视觉主体;
% 底部两层半透明白色波浪;无任何文字。
\newcommand{\makebackcover}{%
    \begingroup
    \let\cleardoublepage\clearpage
    \begin{titlepage}
        \thispagestyle{empty}
        \noindent\hspace*{-\parindent}\begin{tikzpicture}[remember picture,overlay]
            % ---------- 底色:对角渐变 ----------
            \shade[top color=bookblue, bottom color=bookteal!48!bookblue,
                   shading angle=-30]
                (current page.north west) rectangle (current page.south east);

            % ---------- 角部半透明叠层 ----------
            \fill[white,opacity=0.05]
                ([xshift=2.0cm,yshift=1.0cm]current page.north east) circle[radius=7.2cm];
            \fill[white,opacity=0.06]
                ([xshift=-0.5cm,yshift=-1.5cm]current page.north east) circle[radius=3.6cm];

            % ---------- 左下同心圆 ----------
            \foreach \r in {1.0,1.75,2.5,3.25,4.0,4.75}{
                \draw[booklight!40,opacity=0.55,line width=0.55pt]
                    ([xshift=-0.6cm,yshift=-0.4cm]current page.south west) circle[radius=\r cm];
            }
            % ---------- 右下螺旋线 ----------
            \begin{scope}[shift={([xshift=-2.9cm,yshift=3.2cm]current page.south east)}]
                \draw[booklight!45!bookteal,opacity=0.5,line width=0.55pt]
                    plot[domain=0:1260,samples=200,smooth,variable=\t]
                    ({0.0036*\t*cos(\t)},{0.0036*\t*sin(\t)});
            \end{scope}
            % ---------- 左上斜线组 ----------
            \foreach \i in {0,1,2}{
                \draw[bookteal!60!booklight,opacity=0.5,line width=0.6pt]
                    ([xshift={1.6cm+\i*0.45cm}]current page.north west) --
                    ([yshift={-1.6cm-\i*0.45cm}]current page.north west);
            }

            % ---------- 中下部视觉主体:外摆线(玫瑰形) ----------
            \begin{scope}[shift={([yshift=-16.6cm]current page.north)}]
                % 外圈
                \draw[booklight!35,opacity=0.6,line width=0.6pt]
                    (0,0) circle[radius=4.9cm];
                \draw[booklight!22,opacity=0.5,line width=0.5pt]
                    (0,0) circle[radius=5.25cm];
                % 主玫瑰线
                \draw[booklight!55!bookteal,opacity=0.8,line width=0.85pt]
                    plot[domain=0:1080,samples=400,smooth,variable=\t]
                    ({0.335*(8*cos(\t)-5*cos(2.6667*\t))},
                     {0.335*(8*sin(\t)-5*sin(2.6667*\t))});
                % 内层玫瑰线
                \draw[booklight!38!bookteal,opacity=0.6,line width=0.6pt]
                    plot[domain=0:1080,samples=400,smooth,variable=\t]
                    ({0.21*(8*cos(\t)-5*cos(2.6667*\t))},
                     {0.21*(8*sin(\t)-5*sin(2.6667*\t))});
            \end{scope}

            % ---------- 底部半透明波浪 ----------
            \fill[white,opacity=0.06]
                ([yshift=4.6cm]current page.south west)
                .. controls ([xshift=6.5cm,yshift=6.4cm]current page.south west)
                         and ([xshift=-6.0cm,yshift=2.6cm]current page.south east) ..
                ([yshift=5.2cm]current page.south east) --
                (current page.south east) -- (current page.south west) -- cycle;
            \fill[booklight,opacity=0.08]
                ([yshift=2.6cm]current page.south west)
                .. controls ([xshift=5.5cm,yshift=4.3cm]current page.south west)
                         and ([xshift=-6.5cm,yshift=1.0cm]current page.south east) ..
                ([yshift=3.3cm]current page.south east) --
                (current page.south east) -- (current page.south west) -- cycle;
        \end{tikzpicture}
        \mbox{}
    \end{titlepage}
    \endgroup
}
\makeatother

% 页脚右下方常态化水印(署名 + 仓库地址,加粗黑色),已在 main/plain 分页样式中定义,所有页面均生效.

\begin{document}
\mainmatter
\pagestyle{main}
\setcounter{chapter}{4}
\def\BookEnd{\end{document}}
\fi

\chapter{多元微分学与多元积分学}
\begin{ti}
【题1】若$f(x,y)$是在$\mathbb{R}^2\backslash\{0\}$上连续的$n$次齐次函数(注:$n$次齐次函数是指满足$f(t\mathbf{x})=t^nf(\mathbf{x})$的函数),
证明:$\left\lvert f(x,y)\right\rvert\leq C(x^2+y^2)^{\frac{n}{2}}$,其中$C$为正常数
\end{ti}

\textbf{证明}:考虑集合$S=\left\{(x,y)\,|\;x^2+y^2=1\right\}$,则$S$为有界闭集, \\

又$f$为连续函数,$\therefore\left\lvert f\right\rvert $也是连续函数,$\therefore\exists\,C>0,\,(x,y)\in S\text{时},\left\lvert f(x,y)\right\rvert\leq C$, \\

对$\forall\,(x,y)\in\mathbb{R}^2\backslash\{0\},\text{令}r=\sqrt{x^2+y^2},$ 则$f(x,y)=f(r\cdot\dfrac{x}{r},r\cdot\dfrac{y}{r})=r^nf(\dfrac{x}{r},\dfrac{y}{r}),$ \\

又$\because(\dfrac{x}{r},\dfrac{y}{r})\in S,\quad\therefore\left\lvert f(\dfrac{x}{r},\dfrac{y}{r})\right\rvert\leq C$, \\

$\therefore\left\lvert f(x,y)\right\rvert\leq C(x^2+y^2)^{\frac{n}{2}}$.

\begin{ti}
【题2】如果$F(x,y)$在矩形$D=\left\{(x,y)\in\mathbb{R}^2\,|\;\left\lvert x-x_0\right\rvert<a,\left\lvert y-y_0\right\rvert<b\right\}$上连续,
$F(x_0,y_0)=0,$且$F(x,y)$对每一个$x\in(x_0-a,x_0+a)$关于$y$严格单调,求证:在$(x_0,y_0)$某个邻域内,由$F(x,y)=0$可确定唯一的函数$y=f(x)$,且$f(x)$在$x_0$的一个邻域里连续
\end{ti}

\textbf{证明}:不妨设关于$\displaystyle y$严格单调递增.取$\displaystyle(x_0,y_1),(x_0,y_2)\in D$,且$\displaystyle y_1>y_0>y_2$, \\

则$\displaystyle F(x_0,y_1)>0,F(x_0,y_2)<0$.取$\displaystyle0<\alpha<a$使$\displaystyle S=[x_0-\alpha,x_0+\alpha]\times[y_2,y_1]\subset D$, \\

因为$\displaystyle F$在闭矩形$\displaystyle S$上一致连续,令$\displaystyle\varepsilon_1=\dfrac{\min\{|F(x_0,y_1)|,|F(x_0,y_2)|\}}{2}$, \\

可取$\displaystyle\delta_0>0$,使$\displaystyle u,v\in S$且$\displaystyle\left\lVert u-v\right\rVert <\delta_0$时,$\displaystyle|F(u)-F(v)|<\varepsilon_1$, \\

令$\displaystyle\delta=\min\{\delta_0,\alpha\}$,对每个$\displaystyle x\in(x_0-\delta,x_0+\delta)$,有$\displaystyle F(x,y_1)>0,F(x,y_2)<0$, \\

故由连续性和严格单调性,存在唯一的$\displaystyle f(x)\in(y_2,y_1)$使$\displaystyle F(x,f(x))=0$,且$\displaystyle f(x_0)=y_0$, \\

下证$\displaystyle f$连续,任取$\displaystyle x^*\in(x_0-\delta,x_0+\delta)$及$\displaystyle\varepsilon>0$, \\

取$\displaystyle 0<\rho<\min\{\varepsilon,f(x^*)-y_2,y_1-f(x^*)\}$, \\

则由严格单调性,$\displaystyle F(x^*,f(x^*)-\rho)<0<F(x^*,f(x^*)+\rho)$, \\

由$\displaystyle F$的连续性,存在$\displaystyle\eta_0>0$,使$\displaystyle|x-x^*|<\eta_0$时上述两端仍分别为负和正, \\

令$\displaystyle\eta=\min\{\eta_0,\delta-|x^*-x_0|\}>0$,当$\displaystyle|x-x^*|<\eta$时, \\

唯一零点$\displaystyle f(x)$位于$\displaystyle(f(x_*)-\rho,f(x_*)+\rho)$内,故$\displaystyle|f(x)-f(x_*)|<\rho<\varepsilon$, \\

所以$\displaystyle f$在$\displaystyle(x_0-\delta,x_0+\delta)$上连续,命题得证.

\begin{ti}
【题3】$D=\left\{u\in\mathbb{R}^n\,|\;\left\lVert u\right\rVert\leq\frac{1}{2}\right\},$ $f\colon D\to\mathbb{R}\text{为}C^1$映射, \\
$\left\lVert \nabla f(0)\right\rVert=1,\text{且}\left\lVert \nabla f(u)-\nabla f(v)\right\rVert\leq\left\lVert u-v\right\rVert $, \\
证明:存在唯一的点$\xi$,使得$f(\xi)=\displaystyle\max_{D}f$
\end{ti}

\textbf{证明}:$\because D$为有界闭集,$f$为连续映射,从而$f$在$D$上能取到最大值, \\

当$u\in int D\text{时},\left\lVert \nabla f(u)-\nabla f(0)\right\rVert\leq\left\lVert u\right\rVert<\dfrac{1}{2}\Rightarrow\dfrac{1}{2}<\left\lVert \nabla f(u)\right\rVert<\dfrac{3}{2},$ \\

特别地,$\left\lVert \nabla f(u)\right\rVert\neq 0,$从而$f$的最大值只能在$\partial D$上取到, \\

令$F(u)=f(u)-\dfrac{1}{2}\left\lVert u\right\rVert^2,\;\therefore\,\nabla F(u)=\nabla f(u)-u$,同理$F$可在$D$上取到最大值, \\

当$u\in int D\text{时},\left\lVert \nabla F(u)\right\rVert=\left\lVert \nabla f(u)-u\right\rVert\geq\left\lVert \nabla f(u)\right\rVert-\left\lVert u\right\rVert>0$, \\

$\therefore\,\nabla F(u)\neq 0$,即$F$的最大值也只能在$\partial D$上取到, \\

又$\displaystyle\max_{D}F=\max_{\partial D}F=\max_{\partial D}f-\dfrac{1}{8}=\max_{D}f-\dfrac{1}{8}$, \\

$\therefore F\text{与}f$有完全相同的最大值点,因此只需要证明$F$的最大值点唯一即可, \\

对任意$\displaystyle u,v\in D$,由条件与$Cauchy$不等式, \\

$\displaystyle(\nabla F(u)-\nabla F(v))\cdot(u-v)$ \\

$\displaystyle=(\nabla f(u)-\nabla f(v))\cdot(u-v)-\left\lVert u-v\right\rVert^2$ \\

$\displaystyle\leq\left\lVert\nabla f(u)-\nabla f(v)\right\rVert\left\lVert u-v\right\rVert-\left\lVert u-v\right\rVert^2\leq0,$ \\

令$\displaystyle w(t)=(1-t)u+tv,\;\psi(t)=F(w(t))$,若$\displaystyle0\leq s<t\leq1$, \\

则$\displaystyle(\nabla F(w(t))-\nabla F(w(s)))\cdot(w(t)-w(s))\leq0,$ \\

而$\displaystyle w(t)-w(s)=(t-s)(v-u)$,故$\displaystyle\psi'(t)\leq\psi'(s)$, \\

因此$\displaystyle\psi'$单调不增,$\displaystyle\psi$为凹函数,从而$\displaystyle F$在$\displaystyle D$上为凹函数, \\

假设$\displaystyle u^0,v^0$是$\displaystyle F$的两个不同最大值点,记$\displaystyle F(u^0)=F(v^0)=M$, \\

对任意$\displaystyle t\in(0,1)$,由凹性与$\displaystyle M$的最大性, \\

$\displaystyle M\geq F((1-t)u^0+tv^0)\geq(1-t)F(u^0)+tF(v^0)=M,$ \\

故线段开段上每一点都是$\displaystyle F$的最大值点, \\

由于闭球$\displaystyle D$严格凸,$\displaystyle(1-t)u^0+tv^0\in\operatorname{int}D$, \\

由$Fermat$必要条件可知该点满足$\displaystyle\nabla F=0$,与前面已证的$\displaystyle\nabla F\neq0$矛盾, \\

从而$F$只能有唯一的最大值点,即存在唯一的点$\xi$,使得$f(\xi)=\displaystyle\max_{D}f$,得证.

\begin{ti}
【题4】二元函数$f(x,y)$,定义域为$I=[a,b]\times[c,d]$,$f$关于每个变量单调递增,证明:$f\in R(I)$
\end{ti}

\textbf{证明}:取等距分划$\displaystyle a=x_0<\cdots<x_n=b,\quad c=y_0<\cdots<y_n=d$,并令 \\

$\displaystyle R_{ij}=[x_{i-1},x_i]\times[y_{j-1},y_j],\quad\Delta x=\dfrac{b-a}{n},\quad\Delta y=\dfrac{d-c}{n}.$ \\

由单调性,$\displaystyle f$在$\displaystyle R_{ij}$上的振幅为$\displaystyle\omega_{ij}=f(x_i,y_j)-f(x_{i-1},y_{j-1})$ \\

$\displaystyle=[f(x_i,y_j)-f(x_{i-1},y_j)]+[f(x_{i-1},y_j)-f(x_{i-1},y_{j-1})]$, \\

因此沿$\displaystyle i,j$分别求和,可得$\displaystyle S(f,P_n)-s(f,P_n)=\sum_{i,j}\omega_{ij}\Delta x\Delta y$ \\

$\displaystyle=\Delta x\Delta y\left\{\sum_{j=1}^n[f(b,y_j)-f(a,y_j)]+\sum_{i=1}^n[f(x_{i-1},d)-f(x_{i-1},c)]\right\}$ \\

$\displaystyle\leq\frac{2(b-a)(d-c)}{n}[f(b,d)-f(a,c)]\to0.$ \\

故对任意$\displaystyle\varepsilon>0$,可取充分大的$\displaystyle n$,使$\displaystyle S(f,P_n)-s(f,P_n)<\varepsilon$,从而$\displaystyle f\in R(I)$,得证.

\begin{ti}
【题5】$\gamma(t):[a,b]\to \mathbb{R}^2,\gamma(t)=(x(t),y(t))$,且 $x(t)\in C^\alpha[a,b],y(t)\in C^\beta[a,b],0<\alpha,\beta<1, \alpha+\beta>1$, 
证明:$\gamma([a,b])$ 是 $\mathbb{R}^2$ 中零面积集.(注:$x(t)\in C^\alpha[a,b]$指的是$\displaystyle|x(t_1)-x(t_2)|\leq C_x\left\lvert t_1-t_2\right\rvert^{\alpha}$)
\end{ti}

\textbf{证明}:将 $[a,b]$ $n$等分,$a=t_0<t_1<\dots<t_n=b, \Delta t_i=\dfrac{b-a}{n}$, \\

 设$x(t)$在$[t_{i-1},t_i]$上最大值,最小值分别为$\alpha_i, \beta_i$, \\
 
 $\phantom{\text{设}}y(t)$在$[t_{i-1},t_i]$上最大值,最小值分别为$\mu_i, \nu _i$, \\
 
 令$I_i=[\beta_i,\alpha_i]\times[\nu_i,\mu_i]$ ,则 $\displaystyle\gamma([a,b])\subset\bigcup_{i=1}^n I_i$, 取$C=\max\{C_x,C_y\}$, \\
 
 $\because \alpha_i - \beta_i \leq C|t_i - t_{i-1}|^\alpha = \dfrac{C(b-a)^\alpha}{n^\alpha}$, \\

 $\phantom{\because}\mu_i - \nu_i \leq C|t_i - t_{i-1}|^\beta = \dfrac{C(b-a)^\beta}{n^\beta}$, \\
 
 $\displaystyle\therefore \sum_{i=1}^n \sigma(I_i) = \sum_{i=1}^n \dfrac{C^2(b-a)^{\alpha+\beta}}{n^{\alpha+\beta}} = \dfrac{C^2(b-a)^{\alpha+\beta}}{n^{\alpha+\beta-1}}$, \\
 
 $\displaystyle\because \alpha+\beta>1, \quad\therefore \lim_{n\to\infty} \dfrac{C^2(b-a)^{\alpha+\beta}}{n^{\alpha+\beta-1}}=0$, \\

$\displaystyle\therefore \forall\,\varepsilon>0,\;\exists N,\;\text{s.t.}\,\forall\,n\geq N,\;\gamma([a,b])\subset\bigcup_{i=1}^n I_i$,且$\displaystyle\sum_{i=1}^n \sigma(I_i)<\varepsilon$, \\

从而 $\gamma([a,b])$是$\mathbb{R}^2$中零面积集,得证.

\begin{ti}
【题6】设 $f\in C^2(\overline{B_1(0)})$,在 $B_1(0)$ 上 $C^2$,设 $f$ 在 $\partial B_1(0)$ 上最小值为 $a$,
$f(0)=b<a$,考虑映射 $\phi:B_1(0)\to\mathbb{R}^n,\phi(x)=\nabla f(x)$,$\Sigma=\{x\in B_1(0)|\nabla^2 f(x)$半正定$\}$,
证明:$\phi(\Sigma)$ 中包含一个半径为 $a-b$ 的开球
\end{ti}

\textbf{证明}:$\displaystyle\because f(0)=b<a=\min_{\partial B_1(0)} f$,且 $\overline{B_1(0)}$ 是有界闭集, $f\in C^2(\overline{B_1(0)})$, \\

$\therefore\,f$在$\overline{B_1(0)}$ 上有最小值,且最小值点在 $B_1(0)$ 取到, \\

即 $\exists\,x^*\in B_1(0),\;\text{s.t.}\, \nabla f(x^*)=0$,且 $\nabla^2 f(x^*)\geq0,\quad\therefore x^*\in\Sigma$, \\

$\because\,0=\phi(x^*),\quad\therefore 0\in\phi(\Sigma)$, \\

对$\forall v\in B_{a-b}(0)$,考虑函数$g(x)=f(x)-v\cdot x,\;\text{则}g\in C^0(\overline{B_1(0)})$,且$g(0)=f(0)=b$, \\

$x\in\partial B_1(0)$ 时, $g(x)\geq a - v\cdot x \geq a - \left\lVert v\right\rVert\left\lVert x\right\rVert   > a - (a-b)=b$, \\

从而$g(x)$在$\overline{B_1(0)}$上的最小值在$B_1(0)$取到, \\

即$\exists\,y^0\in B_1(0)$,$\;\text{s.t.}\,\nabla g(y^0)=0$且$\nabla^2 g(y^0)\geq0$, \\

又$\nabla g(y^0)=\nabla f(y^0)-v$, $\nabla^2 g(y^0)=\nabla^2 f(y^0)\geq0$, \\

$\therefore\,y^0\in\Sigma$$\quad\therefore\,v=\nabla f(y^0)\in\phi(\Sigma)$, \\

从而$B_{a-b}(0)\subset\phi(\Sigma)$,证毕.

\begin{ti}
【题7】设$D\subset\mathbb{R}^2$为平面区域,$u\in C^2(D)$,满足$\dfrac{\partial^2 u}{\partial x^2}+\dfrac{\partial^2 u}{\partial y^2}=0$,
证明:对任何$\overline{B_r(a)}\subset D$,都有$\displaystyle\int_{\partial B_r(a)}\dfrac{\partial u}{\partial \vec{n}}ds=0$,
其中$\vec{n}$为$\partial B_r(a)$的单位外法向量,$\dfrac{\partial u}{\partial \vec{n} }$为方向导数.
\end{ti}

\textbf{证明}:设$a=(a_1,a_2)$,则$\partial B_r(a)$:$\begin{cases}x=a_1+r\cos\theta\\y=a_2+r\sin\theta\end{cases}\quad\theta\in[0,2\pi]$, \\

$\therefore\,\vec{n}=(\cos\theta,\sin\theta)$, $ds=rd\theta$, \\

又$u\in C^2(D)$,$\quad\therefore\dfrac{\partial u}{\partial \vec{n} }=\nabla u\cdot\vec{n}=\dfrac{\partial u}{\partial x}\cos\theta+\dfrac{\partial u}{\partial y}\sin\theta$,\\

$\therefore\displaystyle\int_{\partial B_r(a)}\dfrac{\partial u}{\partial \vec{n} }ds$
$=\displaystyle\int_0^{2\pi}[\dfrac{\partial u}{\partial x}(a_1+r\cos\theta,a_2+r\sin\theta)\cdot\cos\theta+\dfrac{\partial u}{\partial y}(a_1+r\cos\theta,a_2+r\sin\theta)\sin\theta] rd\theta$ \\

$=\displaystyle\int_0^{2\pi}[\dfrac{\partial u}{\partial x}\cdot(r\cos\theta)+\dfrac{\partial u}{\partial y}\cdot(r\sin\theta)]d\theta$
$=\displaystyle\int_0^{2\pi}[\dfrac{\partial u}{\partial x}\cdot y'(\theta)-\dfrac{\partial u}{\partial y}\cdot x'(\theta)]d\theta$ \\

$=\displaystyle\oint_{\partial B_r(a)}\dfrac{\partial u}{\partial x}dy-\dfrac{\partial u}{\partial y}dx$, \\

设$P(x,y)=-\dfrac{\partial u}{\partial y}$,$\;Q(x,y)=\dfrac{\partial u}{\partial x}$,$\quad\therefore\,P,Q\in C^1(D)$, $\dfrac{\partial Q}{\partial x}=\dfrac{\partial^2 u}{\partial x^2}$,$\;\dfrac{\partial P}{\partial y}=-\dfrac{\partial^2 u}{\partial y^2}$, \\

由$Green$公式,$\displaystyle\oint_{\partial B_r(a)}\dfrac{\partial u}{\partial x}dy-\dfrac{\partial u}{\partial y}dx=\iint_{B_r(a)}(\dfrac{\partial^2 u}{\partial x^2}+\dfrac{\partial^2 u}{\partial y^2})dxdy=0$, \\

$\therefore\displaystyle\int_{\partial B_r(a)}\dfrac{\partial u}{\partial \vec{n} }ds=0$,得证.

\begin{ti}
【题8】$\Omega$为有界$C^1$区域,$\displaystyle u\in C^2(\overline\Omega)$,$\displaystyle\Delta u=u$,且$\displaystyle u|_{\partial\Omega}\equiv 0$,
证明:$u\equiv0$
\end{ti}

\textbf{微分学视角}:假设$u$不恒为$0$,存在$x^0\in int\Omega$,$\;\text{s.t.}\,u(x^0)>0$, \\

且$\displaystyle u(x^0)=\max_{\Omega}u$,则$\nabla u(x^0)=0$,且$\nabla^2 u(x^0)\leq 0$, \\

由于$\Delta u = \dfrac{\partial^2 u}{\partial x^2} + \dfrac{\partial^2 u}{\partial y^2} = \text{tr}(\nabla^2 u)$, \\

$\therefore\,\Delta u(x^0)\leq 0$,与$\Delta u(x^0)=u(x^0)$矛盾!  从而$u \equiv 0$,得证. \\

\textbf{积分学视角}:考虑积分$\displaystyle\int_\Omega u^2 dxdy$,则$\displaystyle\int_\Omega u^2 dxdy = \int_\Omega u\Delta u dxdy$, \\

$\because\,\text{div}(u \nabla u) = \text{div}((u\dfrac{\partial u}{\partial x},u\dfrac{\partial u}{\partial y})) = \dfrac{\partial(u\dfrac{\partial u}{\partial x})}{\partial x}+\dfrac{\partial (u\dfrac{\partial u}{\partial y})}{\partial y}$ \\

$=u\dfrac{\partial^2 u}{\partial x^2} + (\dfrac{\partial u}{\partial x})^2 + u\dfrac{\partial^2 u}{\partial y^2} + (\dfrac{\partial u}{\partial y})^2$ $=u\Delta u+||\nabla u||^2$, \\

由$Gauss$公式,$\displaystyle\int_{\Omega} u\Delta u dxdy = \int_\Omega \text{div}(u \nabla u) dxdy - \int_\Omega ||\nabla u||^2 dxdy$ \\

$\displaystyle=\int_{\partial\Omega} u\nabla u \cdot \vec{n} ds - \int_\Omega ||\nabla u||^2 dxdy$, \\

$\displaystyle=-\int_\Omega ||\nabla u||^2 dxdy \leq 0$,而$\int_\Omega u^2 dxdy \geq 0$, \\

从而$\displaystyle\int_\Omega u^2 dxdy = 0 \Rightarrow u \equiv 0$,得证.

\begin{ti}
【题9】计算$\displaystyle\oint_\Gamma\dfrac{xdy - ydx}{ax^2 + 2cxy + by^2}$,
其中$\Gamma:x^2 + y^2 =1$,方向为逆时针且$ab > c^2$
\end{ti}

\textbf{证明}:分母$\displaystyle ax^2+2cxy+by^2=(x,y)M\begin{pmatrix}x\\y\end{pmatrix}$, \\

其中$\displaystyle M=\begin{pmatrix}a&c\\c&b\end{pmatrix}$为实对称矩阵,且$\displaystyle\det M=ab-c^2>0$, \\

可选$\displaystyle P\in SO(2)$,使$\displaystyle P^TMP=\begin{pmatrix}\lambda_1&0\\0&\lambda_2\end{pmatrix}$, \\

令$\displaystyle\begin{pmatrix}x\\y\end{pmatrix}=P\begin{pmatrix}u\\v\end{pmatrix}$,则单位圆及其逆时针方向保持不变,且$\displaystyle x\,dy-y\,dx=u\,dv-v\,du$, \\

由$\displaystyle ab>c^2$可知$\displaystyle a,b$同号,两个特征值也与$\displaystyle a$同号, \\

令$\displaystyle s=\operatorname{sgn}(a)$,则可写$\displaystyle\lambda_i=s\mu_i$,其中$\displaystyle\mu_i>0$,且$\displaystyle\mu_1\mu_2=\lambda_1\lambda_2=ab-c^2$, \\

取$\displaystyle u=\cos\theta,v=\sin\theta$,则$\displaystyle\oint_\Gamma\dfrac{x\,dy-y\,dx}{ax^2+2cxy+by^2}$ \\

$\displaystyle=s\int_0^{2\pi}\dfrac{d\theta}{\mu_1\cos^2\theta+\mu_2\sin^2\theta}$ $\displaystyle=4s\int_0^{\frac{\pi}{2}}\dfrac{d\theta}{\mu_1\cos^2\theta+\mu_2\sin^2\theta}$, \\

令$\displaystyle t=\tan\theta$,则$\displaystyle4s\int_0^{\frac{\pi}{2}}\dfrac{d\theta}{\mu_1\cos^2\theta+\mu_2\sin^2\theta}=4s\int_0^{+\infty}\dfrac{dt}{\mu_1+\mu_2t^2}=\dfrac{2\pi s}{\sqrt{\mu_1\mu_2}},$ \\

故:$\quad\displaystyle\boxed{\oint_\Gamma\dfrac{x\,dy-y\,dx}{ax^2+2cxy+by^2}=\dfrac{2\pi\operatorname{sgn}(a)}{\sqrt{ab-c^2}}}.$

\begin{ti}
【题10】$\Omega\subset \mathbb{R}^3$为有$C^1$边界的有界区域,$\displaystyle u\in C^2(\overline\Omega,\mathbb{R})$,
$\Delta u = 0$,证明:对任意内点$p \in \Omega$,有:
$\displaystyle u(p)=\dfrac{1}{4\pi}\int_{\partial\Omega}\left(\dfrac{\boldsymbol{r}\cdot\boldsymbol{n}}{r^3} u + \dfrac{1}{r}\dfrac{\partial u}{\partial \boldsymbol{n}} \right)d\sigma$,
其中$p = (x_0, y_0, z_0)$,$\boldsymbol{r} = (x - x_0, y - y_0, z - z_0)$,$\boldsymbol{n}$为$\partial\Omega$单位外法向量
\end{ti}

\textbf{证明}:固定$\displaystyle\varepsilon_0>0$使$\displaystyle\overline{B_{\varepsilon_0}(p)}\subset\Omega$,并取$\displaystyle0<\varepsilon<\varepsilon_0$.因为$\displaystyle\dfrac{\partial u}{\partial \boldsymbol{n}}=\nabla u\cdot\boldsymbol{n}$,令$\displaystyle\boldsymbol{X}=\dfrac{\boldsymbol{r}}{r^3}u+\dfrac{1}{r}\nabla u$,考虑$\displaystyle\Omega\setminus B_\varepsilon(p)$, \\

记$\displaystyle I_1 =\int_{\partial\Omega}(\dfrac{\boldsymbol{r} \cdot \boldsymbol{n}}{r^3} u +\dfrac{1}{r}\dfrac{\partial u}{\partial \boldsymbol{n}}) d\sigma$,
$\displaystyle I_2 = \int_{\partial B_\varepsilon(p)} (\dfrac{\boldsymbol{r} \cdot \boldsymbol{n}}{r^3} u +\dfrac{1}{r}\dfrac{\partial u}{\partial \boldsymbol{n}}) d\sigma$, \\

则由$Gauss$公式,$\displaystyle I_1 + I_2 = \int_{\Omega\backslash B_{\varepsilon}(p)}\text{div}\boldsymbol{X} \, dxdydz$,记$\boldsymbol{X} = (P, Q, R)$, \\

$\therefore\dfrac{\partial P}{\partial x} = \dfrac{1}{r}\dfrac{\partial^2 u}{\partial x^2} + (-\dfrac{1}{r^2})\dfrac{x-x_0}{r}\dfrac{\partial u}{\partial x}+\dfrac{x-x_0}{r^3}\dfrac{\partial u}{\partial x} + u\cdot\dfrac{r^2 - 3(x - x_0)^2}{r^5}$ \\

$=\dfrac{1}{r}\dfrac{\partial^2 u}{\partial x^2} + u\cdot\dfrac{r^2 - 3(x - x_0)^2}{r^5}$, \\

同理$\dfrac{\partial Q}{\partial y} = \dfrac{1}{r} \dfrac{\partial^2 u}{\partial y^2} + u\cdot\dfrac{r^2 - 3(y - y_0)^2}{r^5}$,
$\dfrac{\partial R}{\partial z} = \dfrac{1}{r}\dfrac{\partial^2 u}{\partial z^2} + u\cdot\dfrac{r^2 - 3(z - z_0)^2}{r^5}$, \\

$\therefore\,\text{div} \boldsymbol{X} = \dfrac{1}{r}\Delta u + u\cdot\dfrac{3r^2 - 3r^2}{r^5} = 0$, 
$\quad\therefore I_1 + I_2 = 0$, \\

在$\partial B_\varepsilon(p)$上,诱导定向的$\boldsymbol{n} = -\dfrac{\boldsymbol{r}}{r}$,且$r = \varepsilon$, \\

$\displaystyle\therefore\,I_2 = -\int_{\partial B_\varepsilon(p)}(\dfrac{u}{\varepsilon^2}+\dfrac{\nabla u}{\varepsilon}\cdot\dfrac{\boldsymbol{r}}{r} \,) d\sigma$, \\

令$\displaystyle C=\max_{\overline{B_{\varepsilon_0}(p)}}\left\lVert \nabla u\right\rVert <\infty$.对每个$\displaystyle0<\varepsilon<\varepsilon_0$,有 \\

$\displaystyle\left|\int_{\partial B_\varepsilon(p)}\frac{\nabla u}{\varepsilon}\cdot\frac{\boldsymbol{r}}{r}\,d\sigma\right|\leq\frac{C}{\varepsilon}\,4\pi\varepsilon^2=4\pi C\varepsilon\to0,$ \\

又由$\displaystyle u$在$\displaystyle p$处连续,有$\displaystyle\sup_{|x-p|=\varepsilon}|u(x)-u(p)|\to0$, \\

从而$\displaystyle\left|\frac{1}{\varepsilon^2}\int_{\partial B_\varepsilon(p)}u\,d\sigma-4\pi u(p)\right|\leq4\pi\sup_{|x-p|=\varepsilon}|u(x)-u(p)|\to0,$ \\

$\displaystyle\therefore\lim_{\varepsilon\to0}I_2=-4\pi u(p)$,又$\displaystyle I_1+I_2=0$对每个$\displaystyle0<\varepsilon<\varepsilon_0$成立, \\

$\therefore\,I_1 = 4\pi u(p)$,即$\displaystyle u(p) = \dfrac{1}{4\pi}\int_{\partial \Omega} \left( \dfrac{\boldsymbol{r} \cdot \boldsymbol{n}}{r^3} u + \dfrac{1}{r} \dfrac{\partial u}{\partial \boldsymbol{n}} \right) d\sigma$.



\BookEnd
